% GENERATED FILE. Edit canonical lessons, then run npm run generate:books.
% canonical-source-hash: fnv1a64:601077a68af7134f
% canonical-lessons: TA-C213-alamari, TA-C213-mincaram, TA-C213-kuppai, TA-C213-parpacai, TA-C213-paltulakki

\chapter{Cupboard, Electricity, Rubbish, Toothpaste, Toothbrush}
\label{ch:ta-cupboard-electricity-rubbish-toothpaste-toothbrush}
\hlchaptermodality{213}

\begin{chapteropening}
\textbf{By the end of this chapter:} \emph{I can say a cupboard, electricity, rubbish, toothpaste and a toothbrush.}

Complete the last lesson of chapter 213: I can say a cupboard, electricity, rubbish, toothpaste and a toothbrush.
\end{chapteropening}

\section[\texorpdfstring{\ta{அலமாரி} (alamāri)}{அலமாரி (alamāri)}]{\ta{அலமாரி} (alamāri) — a cupboard}

\label{lesson:TA-C213-alamari}



\begin{quote}
Before the new one: say the Tamil for well-behaved, then the Tamil for ugly.
\end{quote}

\subsection*{You'll want to know: \ta{அலமாரி}}
\textbf{\ta{அலமாரி}} — \emph{alamāri} — \textquotedblleft{}a cupboard\textquotedblright{}.

\begin{grammarlens}[title={What you've built}]
One more word for this chapter.
\end{grammarlens}

\subsection*{Your turn}
\begin{itemize}
\raggedright
  \item \emph{Say it:} \emph{alamāri}
  \item \emph{Say it:} \emph{alamāri}, once more
  \item \emph{Say it:} say \emph{alamāri}
  \item \emph{Recall:} say the Tamil for well-behaved, then the Tamil for ugly, then say \emph{alamāri} again
\end{itemize}

\begin{tcolorbox}[breakable,colback=teal!4,colframe=teal!35!black,title={Before you move on}]
What does \ta{அலமாரி} mean? (\textquotedblleft{}A cupboard\textquotedblright{}.) Say it once more.
\end{tcolorbox}
\section[\texorpdfstring{\ta{மின்சாரம்} (miṉcāram)}{மின்சாரம் (miṉcāram)}]{\ta{மின்சாரம்} (miṉcāram) — electricity}

\label{lesson:TA-C213-mincaram}



\begin{quote}
Before the new one: say the Tamil for ugly, then the Tamil for a cupboard.
\end{quote}

\subsection*{You'll want to know: \ta{மின்சாரம்}}
\textbf{\ta{மின்சாரம்}} — \emph{miṉcāram} — \textquotedblleft{}electricity\textquotedblright{}.

\begin{grammarlens}[title={What you've built}]
One more word for this chapter.
\end{grammarlens}

\subsection*{Your turn}
\begin{itemize}
\raggedright
  \item \emph{Say it:} \emph{miṉcāram}
  \item \emph{Say it:} \emph{miṉcāram}, once more
  \item \emph{Say it:} say \emph{miṉcāram}
  \item \emph{Recall:} say the Tamil for ugly, then the Tamil for a cupboard, then say \emph{miṉcāram} again
\end{itemize}

\begin{tcolorbox}[breakable,colback=teal!4,colframe=teal!35!black,title={Before you move on}]
What does \ta{மின்சாரம்} mean? (\textquotedblleft{}Electricity\textquotedblright{}.) Say it once more.
\end{tcolorbox}
\section[\texorpdfstring{\ta{குப்பை} (kuppai)}{குப்பை (kuppai)}]{\ta{குப்பை} (kuppai) — rubbish, garbage}

\label{lesson:TA-C213-kuppai}



\begin{quote}
Before the new one: say the Tamil for a cupboard, then the Tamil for electricity.
\end{quote}

\subsection*{You'll want to know: \ta{குப்பை}}
\textbf{\ta{குப்பை}} — \emph{kuppai} — \textquotedblleft{}rubbish, garbage\textquotedblright{}.

\begin{grammarlens}[title={What you've built}]
One more word for this chapter.
\end{grammarlens}

\subsection*{Your turn}
\begin{itemize}
\raggedright
  \item \emph{Say it:} \emph{kuppai}
  \item \emph{Say it:} \emph{kuppai}, once more
  \item \emph{Say it:} say \emph{kuppai}
  \item \emph{Recall:} say the Tamil for a cupboard, then the Tamil for electricity, then say \emph{kuppai} again
\end{itemize}

\begin{tcolorbox}[breakable,colback=teal!4,colframe=teal!35!black,title={Before you move on}]
What does \ta{குப்பை} mean? (\textquotedblleft{}Rubbish, garbage\textquotedblright{}.) Say it once more.
\end{tcolorbox}
\section[\texorpdfstring{\ta{பற்பசை} (paṟpacai)}{பற்பசை (paṟpacai)}]{\ta{பற்பசை} (paṟpacai) — toothpaste}

\label{lesson:TA-C213-parpacai}



\begin{quote}
Before the new one: say the Tamil for electricity, then the Tamil for rubbish.
\end{quote}

\subsection*{You'll want to know: \ta{பற்பசை}}
\textbf{\ta{பற்பசை}} — \emph{paṟpacai} — \textquotedblleft{}toothpaste\textquotedblright{}.

\begin{grammarlens}[title={What you've built}]
One more word for this chapter.
\end{grammarlens}

\subsection*{Your turn}
\begin{itemize}
\raggedright
  \item \emph{Say it:} \emph{paṟpacai}
  \item \emph{Say it:} \emph{paṟpacai}, once more
  \item \emph{Say it:} say \emph{paṟpacai}
  \item \emph{Recall:} say the Tamil for electricity, then the Tamil for rubbish, then say \emph{paṟpacai} again
\end{itemize}

\begin{tcolorbox}[breakable,colback=teal!4,colframe=teal!35!black,title={Before you move on}]
What does \ta{பற்பசை} mean? (\textquotedblleft{}Toothpaste\textquotedblright{}.) Say it once more.
\end{tcolorbox}
\section[\texorpdfstring{\ta{பல்துலக்கி} (paltulakki)}{பல்துலக்கி (paltulakki)}]{\ta{பல்துலக்கி} (paltulakki) — a toothbrush}

\label{lesson:TA-C213-paltulakki}



\begin{quote}
Before the new one: say the Tamil for rubbish, then the Tamil for toothpaste.
\end{quote}

\subsection*{You'll want to know: \ta{பல்துலக்கி}}
\textbf{\ta{பல்துலக்கி}} — \emph{paltulakki} — \textquotedblleft{}a toothbrush\textquotedblright{}.

\begin{grammarlens}[title={What you've built}]
One more word for this chapter.
\end{grammarlens}

\subsection*{Your turn}
\begin{itemize}
\raggedright
  \item \emph{Say it:} \emph{paltulakki}
  \item \emph{Say it:} \emph{paltulakki}, once more
  \item \emph{Say it:} say \emph{paltulakki}
  \item \emph{Recall:} say the Tamil for rubbish, then the Tamil for toothpaste, then say \emph{paltulakki} again
\end{itemize}

\begin{tcolorbox}[breakable,colback=teal!4,colframe=teal!35!black,title={Before you move on}]
What does \ta{பல்துலக்கி} mean? (\textquotedblleft{}A toothbrush\textquotedblright{}.) Say it once more.
\end{tcolorbox}
